\chapter{LITERATURE REVIEW}

\label{ch:literature}

\section{Degree Audit and Course Planning Systems}

Institutional planning tools fall into two families. Degree-audit systems such as DegreeWorks
apply a program's requirement rules to a transcript and report what remains; planning systems such
as Stellic and Coursicle arrange courses into terms and flag conflicts. Both families assume that
the institution has already encoded its rules authoritatively and that offerings, meeting times and
rooms are known. That assumption is reasonable for a university with a mature catalog system and
unreasonable for a student prototype, where the program rules are published as a web page and the
schedule is not published in machine-readable form at all.

Tam and Lam \cite{tam2019} describe a web-based course planning and monitoring system for a Hong
Kong university and show that a domain-specific planner is feasible without an institutional
integration project. Their system checks prerequisites and lays out semesters but does not
represent shared courses across programs, and it relies on institution-supplied rule data. The
distinction matters here: the present work has to {\em derive} and then {\em distrust} its program
data, and it records that distrust in the software rather than in a caveat.

\section{Prerequisite Structures}

Treating a curriculum as a directed acyclic graph over courses is the standard formalisation.
Parhami \cite{parhami2015} analyses curriculum design with graph-theoretic tools, including
topological ordering and critical-path reasoning, and shows how prerequisite structure constrains
the feasible orderings of a program. Wang et al.\ \cite{wang2020} visualise prerequisite graphs
for planning and report that the visual form helps students reason about chains they otherwise
trace by hand.

Two lessons were taken from this literature and both changed the implementation. The first is that
the graph is a {\em planning} structure, not a {\em validity} structure: a topological order says
what is possible, not what the institution will offer in a given term. The second is that the graph
is program-specific. Parhami's analysis is of one program; when a shared course appears in two
programs with different neighbours, a single global edge set is simply the union of two graphs that
the institutions never declared jointly. The artifact therefore keeps courses global and edges
per program, and refuses to fall back to a global graph when a program context is present.

\section{Course Sequencing and Timetabling}

The semester-planning problem is a constrained sequencing problem with a long timetabling
literature. Schaerf \cite{schaerf1999} surveys automated timetabling and separates
{\em timetabling} proper, the assignment of events to times and places under hard and soft
constraints, from the coarser problem of sequencing a student's courses. The present work deliberately solves
only the coarser problem for students, and states that no offering or room data exists to justify
a timetable. For the administrative simulation it solves a capacity-constrained allocation problem
in the simpler form of a scarcity-ordered greedy round, and reports the result as a simulation
rather than a timetable.

\section{Estimating Difficulty from Sparse Ratings}

Course difficulty is used throughout the system as a planning signal, and ratings are sparse: a
new course may have no votes at all. Reporting the raw mean of one or two votes would make the
recommendation engine unstable and would present sampling noise as knowledge. The standard remedy
is shrinkage: an estimate is pulled toward a prior in proportion to how little data supports it,
with the weight on the prior decreasing as the sample grows. This is the empirical-Bayes
construction described by Gelman et al.\ \cite{gelman2013} for hierarchical models. The artifact
applies the elementary form of that estimator -- a weighted mean between the course's own votes and
a program-level prior whose strength is a fixed pseudo-count -- and always displays the number of
votes next to the estimate so a reader can judge how much of the number came from data.

\section{Transactional Integrity in Web Applications}

Several behaviours in the artifact are correctness properties rather than features: a completion
must not be recorded when a prerequisite is unmet; a retry of a lost request must not create a
second grade attempt; two administrators editing the same resource configuration must not silently
overwrite one another. Kleppmann \cite{kleppmann2017} sets out the failure modes these properties
guard against, and Cahill et al.\ \cite{cahill2009} describe serializable snapshot isolation, the
mechanism PostgreSQL exposes as {\tt SERIALIZABLE}. The implementation uses that isolation level
for every multi-row decision, holds authorization rows with explicit locks so that a role change
cannot commit mid-decision, and treats a serialization failure as a bounded retry rather than an
error surfaced to the user. The expected-scope precondition carried by each mutation is a form of
optimistic concurrency control in the same tradition.

\section{Accessibility of Planning Interfaces}

A planning interface is dense: grids of small cards, tables of numbers, and controls that change
what is shown. Web Content Accessibility Guidelines \cite{wcag22} require, among other things, that
controls have accessible names that contain their visible labels, that colour is never the only
carrier of meaning, that focus remains visible and does not disappear when a control becomes
disabled, and that zoom controls behave as described at their limits. Automated rule engines such
as axe-core \cite{axecore} detect a subset of these problems, and only in a real rendering engine,
because a DOM emulator computes no colours and does not move focus the way a browser does. The
artifact was therefore audited in headless Chromium with computed styles rather than only in a unit
test environment, and the audit produced both defects and a deliberate non-fix, both reported in
Chapter~\ref{ch:results}.

\section{Gap Analysis}

The literature supplies the graph formalisation, the sequencing problem, the shrinkage estimator
and the transactional vocabulary that this work uses. What it does not supply, and what the artifact
contributes, is a system that carries {\em provenance} through all of those layers at once: a
program model that distinguishes a reference import from a verified cohort, an account scope that
is enforced in the same transaction as the write, a rating estimate that reports its own
confidence, a saved-plan editor whose retries are provably idempotent, and an allocation simulation
whose contracts are frozen and replayed rather than re-interpreted. The gap is not algorithmic;
it is the engineering discipline of keeping an unverified planning tool honest about what it knows.
